% Part: first-order-logic
% Chapter: completeness
% Section: lindenbaums-lemma

\documentclass[../../../include/open-logic-section]{subfiles}

\begin{document}

\iftag{FOL}
      {\olfileid{fol}{com}{lin}}
      {\olfileid{pl}{com}{lin}}

\olsection{Het lemma van Lindenbaum}

\begin{explain}
We bewijzen nu dat elke consistente verzameling !!{sentence}s in een
grotere verzameling past die niet alleen consistent, maar ook
!!{complete} is. We bouwen die grotere verzameling stap voor stap. Bij
elke stap voegen we één !!{sentence} toe en zorgen we dat de verzameling
consistent blijft. We doen dit zo dat voor elke $!A$ op een bepaald
moment $!A$ of $\lnot !A$ wordt toegevoegd. Daarna nemen we de vereniging
van alle stappen. Die bevat dan $!A$ of haar ontkenning~$\lnot !A$ en is
dus volledig. Zij is ook consistent, want bij geen enkele stap hebben we
een inconsistentie ingevoerd.
\end{explain}

\begin{lem}[Lemma van Lindenbaum]
\ollabel{lem:lindenbaum} Elke consistente verzameling~$\Gamma$ in een
taal~$\Lang{L}$ past in een grotere verzameling~$\Gamma^*$ die zowel
!!{complete} als consistent is.
\end{lem}

\begin{proof}
Neem een consistente verzameling $\Gamma$. Zet $!A_0$, $!A_1$,
\dots{} in één lijst met alle !!{sentence}s van~$\Lang L$. Begin met
$\Gamma_0 = \Gamma$ en definieer
\[
\Gamma_{n+1} =
\begin{cases}
\Gamma_n \cup \{ !A_n \} & \textrm{als $\Gamma_n \cup \{!A_n\}$
  consistent is;} \\
\Gamma_n \cup \{ \lnot !A_n \} & \textrm{anders.}
\end{cases}
\]
Neem ten slotte $\Gamma^* = \bigcup_{n \geq 0} \Gamma_n$.

Elke $\Gamma_n$ is consistent. Om te beginnen is $\Gamma_0$ per
definitie consistent. Als $\Gamma_{n+1} = \Gamma_n \cup \{!A_n\}$,
hebben we deze mogelijkheid gekozen omdat de nieuwe verzameling
consistent is. Is zij niet consistent, dan nemen we
$\Gamma_{n+1} = \Gamma_n \cup \{\lnot !A_n\}$. We moeten nog nagaan dat
$\Gamma_n \cup \{\lnot !A_n\}$ dan wel consistent is. Stel dat ook deze
verzameling inconsistent is. Dan zijn \emph{zowel}
$\Gamma_n \cup \{!A_n\}$ als $\Gamma_n \cup \{\lnot !A_n\}$
inconsistent. Volgens
\iftag{FOL}{%
  \tagrefs{prfAX/{fol:axd:prv:prop:provability-exhaustive},
    prfSC/{fol:seq:prv:prop:provability-exhaustive},
    prfND/{fol:ntd:prv:prop:provability-exhaustive},
    prfTab/{fol:tab:prv:prop:provability-exhaustive}}}{%
  \tagrefs{prfAX/{pl:axd:prv:prop:provability-exhaustive},
    prfSC/{pl:seq:prv:prop:provability-exhaustive},
    prfND/{pl:ntd:prv:prop:provability-exhaustive},
    prfTab/{pl:tab:prv:prop:provability-exhaustive}}}
zou $\Gamma_n$ zelf dan inconsistent zijn. Dat botst met de
inductieaanname.

Voor elke~$n$ en elke $i < n$ geldt $\Gamma_i \subseteq \Gamma_n$. Dat
bewijzen we met een eenvoudige inductie naar~$n$. Als $n=0$, bestaat er
geen $i < 0$. De bewering klopt dan vanzelf. Neem voor de inductiestap
aan dat zij klopt voor~$n$. We laten zien: als $i < n+1$, dan
$\Gamma_i \subseteq \Gamma_{n+1}$. Onze constructie geeft
$\Gamma_{n+1} = \Gamma_n \cup \{!A_n\}$ of $= \Gamma_n \cup \{\lnot
!A_n\}$. In beide gevallen is $\Gamma_n \subseteq \Gamma_{n+1}$. Stel
nu dat $i < n+1$. Dan geldt $\Gamma_i \subseteq \Gamma_n$ volgens de
inductieaanname als $i < n$, of gebruiken we het duidelijke feit
$\Gamma_n \subseteq \Gamma_n$ als $i = n$. In beide gevallen
krijgen we $\Gamma_i \subseteq \Gamma_{n+1}$ door de transitiviteit
van~$\subseteq$.

Nu kunnen we bewijzen dat $\Gamma^*$ consistent is. Neem een eindige
verzameling $\Gamma' \subseteq \Gamma^*$. Elk $!B \in \Gamma'$ zit ook
in~$\Gamma_i$ voor een bepaalde~$i$. Kies als $n$ de grootste van deze
indices. Omdat $\Gamma_i \subseteq \Gamma_n$ als $i \le n$, geldt
voor elk $!B \in \Gamma'$ ook $\in \Gamma_n$. Anders gezegd:
$\Gamma' \subseteq \Gamma_n$. De verzameling $\Gamma_n$ is consistent.
Dus is elke eindige deelverzameling $\Gamma' \subseteq \Gamma^*$
consistent. Volgens \iftag{FOL}{%
  \tagrefs{prfAX/{fol:axd:ptn:prop:proves-compact},
    prfSC/{fol:seq:ptn:prop:proves-compact},
    prfND/{fol:ntd:ptn:prop:proves-compact},
    prfTab/{fol:tab:ptn:prop:proves-compact}}}{%
  \tagrefs{prfAX/{pl:axd:ptn:prop:proves-compact},
    prfSC/{pl:seq:ptn:prop:proves-compact},
    prfND/{pl:ntd:ptn:prop:proves-compact},
    prfTab/{pl:tab:ptn:prop:proves-compact}}} is $\Gamma^*$ daarom
consistent.

Elke !!{sentence} van $\Frm[L]$ staat in de lijst waarmee we
~$\Gamma^*$ hebben gedefinieerd. Stel dat $!A_n \notin \Gamma^*$. Dat
kan alleen komen doordat $\Gamma_n \cup \{!A_n\}$ inconsistent was. In
dat geval geldt $\lnot !A_n \in \Gamma^*$. De verzameling $\Gamma^*$ is
dus !!{complete}.
\end{proof}

\end{document}
